% Source: word_embeddings.tex
\begin{notechapter}
% ---- chapter-local definitions ----
\definecolor{codegreen}{rgb}{0,0.6,0}
\definecolor{codegray}{rgb}{0.5,0.5,0.5}
\definecolor{codepurple}{rgb}{0.58,0,0.82}
\definecolor{backcolour}{rgb}{0.95,0.95,0.92}
\lstdefinestyle{mystyle}{
    backgroundcolor=\color{backcolour},
    commentstyle=\color{codegreen},
    keywordstyle=\color{magenta},
    numberstyle=\tiny\color{codegray},
    stringstyle=\color{codepurple},
    basicstyle=\ttfamily\footnotesize,
    breakatwhitespace=false,
    breaklines=true,
    captionpos=b,
    keepspaces=true,
    numbers=left,
    numbersep=5pt,
    showspaces=false,
    showstringspaces=false,
    showtabs=false,
    tabsize=2,
    frame=single,
    rulecolor=\color{black},
}
\lstset{style=mystyle}
% -----------------------------------
\chapter{Word Embeddings}\label{ch:word-embeddings}

\section{Introduction}
Word embeddings are dense vector representations of words that 
capture semantic relationships. They are a cornerstone of modern 
natural language processing (NLP) and are used in tasks like 
machine translation, sentiment analysis, and text classification. 
This tutorial explains what word embeddings are, how they are 
obtained, and how they are used, with a focus on practical 
implementation details.

\section{What Are Word Embeddings?}
Word embeddings map words to dense vectors in a continuous vector 
space. Words with similar meanings are mapped to nearby points, 
enabling machines to understand semantic relationships. For 
example:
\begin{itemize}
    \item ``king'' - ``man'' + ``woman'' $\approx$ ``queen''
    \item ``Paris'' - ``France'' + ``Italy'' $\approx$ ``Rome''
\end{itemize}

\section{How Are Word Embeddings Obtained?}
Word embeddings are learned from large text corpora using 
algorithms like Word2Vec, GloVe, and FastText. These algorithms 
analyze word co-occurrence patterns to capture semantic 
relationships.

\subsection{Word2Vec}
Word2Vec is a shallow neural network with two architectures:
\begin{itemize}
    \item \textbf{Skip-Gram}: Predicts context words given a 
    target word.
    \item \textbf{CBOW (Continuous Bag of Words)}: Predicts a 
    target word given its context words.
\end{itemize}

\subsubsection{Skip-Gram Architecture}
The Skip-Gram model takes a target word as input and predicts its 
context words. For example, given the sentence ``The cat sat on 
the mat,'' if the target word is ``sat,'' the model tries to 
predict ``cat'' and ``on.''

\begin{itemize}
    \item \textbf{Input}: One-hot encoded vector of the target 
    word.
    \item \textbf{Hidden Layer}: A weight matrix (the embedding 
    matrix) maps the input to a dense vector.
    \item \textbf{Output}: A softmax layer predicts the 
    probability distribution of context words.
\end{itemize}

The embedding matrix is learned by minimizing the prediction 
error using gradient descent.

\subsubsection{CBOW Architecture}
The CBOW model takes context words as input and predicts the 
target word. For example, given the context words ``cat'' and 
``on,'' the model tries to predict ``sat.''

\begin{itemize}
    \item \textbf{Input}: Average of one-hot encoded vectors of 
    context words.
    \item \textbf{Hidden Layer}: A weight matrix maps the input 
    to a dense vector.
    \item \textbf{Output}: A softmax layer predicts the 
    probability distribution of the target word.
\end{itemize}

\subsubsection{Pretext Task for Learning Embeddings}
The prediction task in Word2Vec (whether Skip-Gram or CBOW) is a
\textbf{pretext task}. The primary goal is not to perform the
prediction itself but to use it as a mechanism to learn meaningful
word embeddings. By training the model to predict words based on
their context (or vice versa), the model learns to place words with
similar meanings or usage patterns close to each other in the vector space. For example:
\begin{itemize}
    \item Words like "king" and "queen" will have similar embeddings
    because they often appear in similar contexts.
    \item The embeddings capture semantic relationships, such as
    analogies (e.g., "king" - "man" + "woman" $\approx$ "queen"). 
\end{itemize}
This approach is based on the \textbf{distributional hypothesis}, which
states that words appearing in similar contexts tend to have similar
meanings.

\subsection{GloVe}
GloVe (Global Vectors for Word Representation) learns embeddings 
by factorizing a word co-occurrence matrix. The key idea is to 
capture global statistics of word co-occurrences.

\subsubsection{Co-Occurrence Matrix}
The co-occurrence matrix \( X \) stores how often words appear 
together in a context window. For example, \( X_{ij} \) is the 
number of times word \( j \) appears in the context of word \( i \).

\subsubsection{Objective Function}
GloVe minimizes the following objective function:
\[
J = \sum_{i,j=1}^V f(X_{ij}) (w_i^T \tilde{w}_j + b_i + \tilde{b}_j - \log X_{ij})^2
\]
where:
\begin{itemize}
    \item \( w_i \) and \( \tilde{w}_j \) are word vectors.
    \item \( b_i \) and \( \tilde{b}_j \) are bias terms.
    \item \( f(X_{ij}) \) is a weighting function to reduce the 
    influence of rare co-occurrences.
\end{itemize}

\subsection{FastText}
FastText extends Word2Vec by representing words as bags of 
character n-grams. This allows it to generate embeddings for rare 
and out-of-vocabulary words.

\subsubsection{Subword Information}
Each word is represented as a sum of its character n-grams. For 
example, the word ``apple'' might be broken into n-grams like 
``app,'' ``ppl,'' ``ple.''

\subsubsection{Training}
FastText uses the same architectures as Word2Vec (Skip-Gram or 
CBOW) but learns embeddings for both words and their subword 
components.

\subsection{Are Word Embeddings Extracted from Deep Networks?}
The answer depends on the type of word embedding model:

\subsubsection{Traditional Embeddings (Word2Vec, GloVe, FastText)}
\begin{itemize}
    \item \textbf{Not Extracted from Deep Networks}: These models 
    use shallow neural networks or matrix factorization.
    \item In Word2Vec, the embeddings are the weights of the 
    hidden layer in a shallow network.
    \item In GloVe, the embeddings are obtained by factorizing a 
    co-occurrence matrix.
    \item In FastText, the embeddings are learned using a shallow 
    network with subword information.
\end{itemize}

\subsubsection{Contextualized Embeddings (BERT, GPT, ELMo)}
\begin{itemize}
    \item \textbf{Extracted from Deep Networks}: These models use 
    deep neural networks (e.g., transformers, LSTMs) to generate 
    embeddings.
    \item In BERT, embeddings are extracted from the final hidden 
    states of a transformer encoder.
    \item In GPT, embeddings are extracted from the final hidden 
    states of a transformer decoder.
    \item In ELMo, embeddings are extracted from the hidden 
    states of a bidirectional LSTM.
\end{itemize}

\section{How Are Word Embeddings Used?}
Word embeddings are used in various NLP tasks:
\begin{itemize}
    \item \textbf{Text Representation}: Words, sentences, or 
    documents are represented as vectors.
    \item \textbf{Semantic Similarity}: Cosine similarity measures 
    how close two words are in vector space.
    \item \textbf{Machine Translation}: Embeddings capture 
    cross-lingual semantic relationships.
    \item \textbf{Sentiment Analysis}: Words are represented as 
    vectors for classification tasks.
    \item \textbf{Named Entity Recognition (NER)}: Embeddings help 
    identify entities like names and dates.
\end{itemize}

\section{Implementing Word Embeddings}
Word embeddings are typically implemented using an \textbf{index 
mapping} (like a dictionary) and an \textbf{embedding matrix}.

\subsection{Index Mapping (Dictionary)}
A dictionary maps each word to a unique index:
\begin{lstlisting}[language=Python]
index_mapping = {
    "king": 0,
    "queen": 1,
    "man": 2,
    "woman": 3
}
\end{lstlisting}

\subsection{Embedding Matrix}
The embedding matrix stores the actual word embeddings. Each row 
corresponds to a word's index:
\begin{lstlisting}[language=Python]
embedding_matrix = [
    [0.7, 0.8, 0.9],  # Index 0: "king"
    [0.6, 0.7, 0.8],  # Index 1: "queen"
    [0.5, 0.6, 0.7],  # Index 2: "man"
    [0.4, 0.5, 0.6]   # Index 3: "woman"
]
\end{lstlisting}

\subsection{Retrieving an Embedding}
To retrieve an embedding:
\begin{enumerate}
    \item Look up the word's index in the dictionary.
    \item Use the index to access the corresponding row in the 
    embedding matrix.
\end{enumerate}
Example:
\begin{lstlisting}[language=Python]
word = "king"
index = index_mapping[word]  # Get index (e.g., 0)
embedding = embedding_matrix[index]  # Get embedding
print(embedding)  # Output: [0.7, 0.8, 0.9]
\end{lstlisting}

\subsection{Why Not Use a Python Dictionary?}
While a dictionary could store embeddings directly, it is less 
efficient for large vocabularies. The embedding matrix approach 
is:
\begin{itemize}
    \item \textbf{Memory-Efficient}: Stores only vectors, not 
    keys.
    \item \textbf{Fast}: Accessing rows in a matrix is faster than 
    dictionary lookups.
    \item \textbf{Compatible}: Works well with machine learning 
    frameworks like TensorFlow and PyTorch.
\end{itemize}

\section{Loading Pretrained Embeddings}
Pretrained embeddings (e.g., Word2Vec, GloVe) can be loaded using 
libraries like \texttt{gensim}, \texttt{spaCy}, or 
\texttt{transformers}.

\subsection{Using Gensim}
\begin{lstlisting}[language=Python]
from gensim.models import KeyedVectors

# Load the pretrained Word2Vec model
model = KeyedVectors.load_word2vec_format(
    'path/to/word2vec.bin', binary=True)

# Get the embedding for a word
word = "king"
embedding = model[word]
print(embedding)
\end{lstlisting}

\subsection{Using spaCy}
\begin{lstlisting}[language=Python]
import spacy

# Load the spaCy model
nlp = spacy.load('en_core_web_sm')

# Get the embedding for a word
word = "king"
embedding = nlp(word).vector
print(embedding)
\end{lstlisting}

\section{Key Takeaways}
\begin{itemize}
    \item Word embeddings are dense vector representations of 
    words.
    \item They are learned using algorithms like Word2Vec, GloVe, 
    or FastText.
    \item In Word2Vec, the prediction task (Skip-Gram or CBOW) is a 
    pretext task used to learn meaningful word embeddings.
    \item Pretrained embeddings can be loaded and used for various 
    NLP tasks.
    \item The mapping from words to embeddings is implemented 
    using an index dictionary and an embedding matrix for 
    efficiency.
    \item Contextualized embeddings (e.g., BERT) generate dynamic 
    representations based on context.
\end{itemize}

\section{Conclusion}
Word embeddings are a powerful tool for representing words in NLP 
tasks. By understanding how they are obtained, implemented, and 
used, you can leverage them to improve the performance of your NLP 
models. Whether you're working with traditional embeddings like 
Word2Vec or modern contextualized embeddings like BERT, the 
principles remain the same: capture semantic relationships and 
represent words in a way that machines can understand.
\end{notechapter}
